\section{Finite algebraic sections: a concrete repair of the Diophantine interface}
\label{sec:finite-sections}

The repository's general Siegel route already separates two assertions:
every supplied finite tuple has a cover by a bounded number of proper
linear subspaces, and each such subspace contains only a bounded number
of curve points. Its unresolved geometric input is a zero-count theorem
for functions with controlled poles. We now prove an elementary substitute
for a substantial explicit family. The resulting theorem supplies the
second assertion without a theory of divisors or Riemann--Roch spaces.
The quantitative approximation theorem supplying the first assertion
remains a separate, unproved input \cite{repo-siegel}.

All rational numbers in this section have canonical signed-numerator,
positive-denominator codes. Equality and order are decidable integer
tests. A polynomial is a finite coefficient array, with trailing zero
coefficients removed when its degree is computed. A point is a pair of
such rational codes. These choices make distinctness, substitution,
matrix rank, and list membership primitive recursive predicates.

\subsection{A finite polynomial zero-count lemma}

\begin{lemma}[Zeros in a supplied finite list]
\label{lem:section-root-count}
Let $h\in\Q[X]$ be nonzero of degree $D$. A finite list of distinct
rational roots of $h$ has length at most $D$. This statement has a proof
in $\iPRA$, uniformly in the coefficient and root lists.
\end{lemma}
\begin{proof}
We first give the division operation, including its bound. If
$h(X)=a_0+\cdots+a_DX^D$ with $D>0$, and $z\in\Q$, define
$q_{D-1}=a_D$ and, going downwards, $q_{j-1}=a_j+zq_j$ for
$1\le j<D$. Put $q(X)=\sum_{j<D}q_jX^j$. Comparing coefficients gives
\[
 h(X)=(X-z)q(X)+a_0+zq_0,
 \qquad a_0+zq_0=h(z).
\]
The computation uses exactly $D$ coefficient updates. If $h(z)=0$,
the remainder vanishes and $q$ has degree $D-1$ with the same nonzero
leading coefficient. If $w\ne z$ is another root, evaluation gives
$0=(w-z)q(w)$; cancellation in $\Q$ gives $q(w)=0$.

Suppose a supplied root list has at least $D+1$ entries, and retain its
first $D+1$ entries $z_0,\ldots,z_D$. Divide successively by
$X-z_0,\ldots,X-z_{D-1}$. After $j$ divisions the quotient has degree
$D-j$, has the original nonzero leading coefficient, and vanishes at
each of the remaining listed roots. At $j=D$ it is a nonzero constant
which vanishes at $z_D$, a contradiction. If $D=0$, the first root
already gives this contradiction.

This is finite induction, not induction on an assertion about all
rational roots. The remaining-root condition is a bounded conjunction
over the fixed input list. Its characteristic Boolean, the polynomial
coefficient identity, and the leading-coefficient equality form a
quantifier-free invariant. All divisions just described are by a monic
linear factor and require no search. The proof therefore uses only the
allowed induction instances.
\end{proof}

\subsection{An explicit annihilator of a proper span}

We describe the required linear algebra so that the word ``proper''
does not conceal an infinite-dimensional selection principle. For a
$t$ by $n$ rational matrix $W$, define its rank to be the largest size
of a nonzero minor, with rank zero when every entry is zero. This is
a bounded finite test. A proper span in $\Q^n$ can be represented by
such a matrix with rank less than $n$.

\begin{lemma}[A bounded, fraction-free annihilator]
\label{lem:section-annihilator}
For $n\ge1$, from a rational $t$ by $n$ matrix $W$ of rank less than
$n$, a primitive recursive operation constructs a nonzero integer
vector $\lambda\in\mathbb Z^n$ satisfying $W\lambda=0$.
Its correctness has a quantifier-free finite-arithmetic proof.
\end{lemma}
\begin{proof}
Clear all denominators using their positive product $D_0$ and put
$V=D_0W$. This does not change which minors vanish. By bounded
enumeration choose a nonzero minor of maximum size $k<n$, with row
set $I$ and column set $J$; let its matrix be $A$. For $k=0$ use the
empty minor, whose determinant is one. Choose the least column
$c\notin J$, and let $b$ be column $c$ restricted to the rows in $I$.
Define
\[
 \lambda_c=\det A,\qquad
 \lambda_J=-\operatorname{adj}(A)b,\qquad
 \lambda_j=0\quad(j\notin J\cup\{c\}).
\]
When $k=0$ this simply means $\lambda_c=1$ and all other entries zero.
In every case $\lambda$ is nonzero.

The selected rows annihilate $\lambda$ because cofactor expansion gives
$A\operatorname{adj}(A)=(\det A)I_k$. For any other row $v$ of $V$,
the block-determinant identity is
\[
 \det\begin{pmatrix} A&b\\v_J&v_c\end{pmatrix}
       =v_c\det A-v_J\operatorname{adj}(A)b=v\lambda.
\]
This identity follows by expansion along the last row; equivalently,
multiply the usual row-subtraction computation by $\det A$ and use
the adjugate identity. The displayed determinant is a minor of size
$k+1$, up to row and column ordering. It vanishes by maximality of
$k$. Thus every row annihilates $\lambda$. For $k=0$, maximality says
that all matrix entries are zero, which gives the same conclusion.
Since $D_0>0$, $V\lambda=0$ implies $W\lambda=0$.

There are at most $2^{t+n}$ row-mask/column-mask pairs to inspect.
Each determinant is the signed sum over permutations, enumerated by
a bounded loop (even a scan of all tuples below $n^n$, retaining the
permutations, suffices). The cofactor identities follow by finite
expansion: terms with repeated columns cancel in pairs, and the
remaining terms are exactly the determinant terms. The block identity
is the same finite expansion. All matrix entries and identities are
finite arrays of rational or integer codes.

For a numerical size bound, if all input numerators in absolute value
and denominators are at most $H\ge1$, then
$D_0\le H^{tn}$ and every entry of $V$ has absolute value at most
$C=H^{tn+1}$. Every coordinate constructed above has absolute value
at most $n!C^n$: a $k$-minor has bound $k!C^k$, and the adjugate
product has that same bound. This is a uniform primitive recursive
majorant. No arbitrary rational kernel vector is selected.
\end{proof}

\subsection{A zero-count theorem for superelliptic curves}

Fix integers $r\ge1$ and $s\ge0$, and a nonzero rational polynomial
$f$ of degree $m$ satisfying
\[
                         m>rs.
\]
Consider the affine point predicate and feature vector
\[
 \mathcal C_{f,r}(x,y)\;:\Longleftrightarrow\;y^r=f(x),
 \qquad
 \Phi_s(x,y)=(1,x,\ldots,x^s,y)\in\Q^{s+2}.
\]
No smoothness, irreducibility, or genus hypothesis is needed for the
following finite assertion.

\begin{theorem}[Uniform finite section bound]
\label{thm:superelliptic-section}
For every nonzero rational vector
$\lambda=(\lambda_0,\ldots,\lambda_s,a)$, a finite list of distinct
points satisfying
\[
             y^r=f(x),\qquad \lambda\cdot\Phi_s(x,y)=0
\]
has at most $m$ entries. Consequently the same bound holds for points
whose feature vectors lie in any supplied proper rational linear
subspace. Both statements have proofs in $\iPRA$.
\end{theorem}
\begin{proof}
Put $P(X)=\sum_{j=0}^s\lambda_jX^j$. The second equation is
$P(x)+ay=0$. The rational test $a=0$ gives two decidable cases.

If $a\ne0$, every such point has $y=-P(x)/a$. Substituting this
into the first equation gives
\[
             H(x)=0,\qquad H(X)=(-P(X))^r-a^rf(X).
\]
The first summand has degree at most $rs<m$, whereas the second has
degree exactly $m$ and nonzero leading coefficient. Hence $H$ is
nonzero of degree exactly $m$. Distinct points in this branch have
distinct $x$ coordinates, because $x$ determines $y$. Applying
Lemma~\ref{lem:section-root-count} to those coordinates gives the
bound $m$.

If $a=0$, the vector's nonzeroness says that $P$ is a nonzero
polynomial of degree at most $s$. Group the input points by equal
$x$ coordinates, assigning each point to the first earlier equal
coordinate, or starting a new group. There are at most $s$ groups,
by Lemma~\ref{lem:section-root-count}. In each group, the distinct
$y$ coordinates are roots of $Y^r-f(x)$, a polynomial of degree
exactly $r$. Thus each group has at most $r$ entries. The entire
list has at most $rs<m$ entries. This also handles $s=0$: a
nonzero constant $P$ has no roots, so there are no groups.

For a proper subspace supplied by a spanning matrix of rank less
than $s+2$, Lemma~\ref{lem:section-annihilator} supplies a nonzero
annihilator. Every vector in its span satisfies the resulting
linear equation. Thus the already proved hyperplane bound applies.

The polynomial powers use finite convolution of coefficient lists.
The degree statements are leading-coefficient statements; the
condition $m>rs$ makes the relevant coefficient survive without
cancellation. Grouping a list of length $t$ requires at most $t^2$
equality tests. The group-size assertion is a bounded conjunction,
and its sum is a primitive recursive fold. Together with the two
lemmas, this gives an explicit reduction of the whole proof to
quantifier-free induction on bounded arrays and loops.
\end{proof}

\begin{remark}[The constant is sharp]
For any $r,s,m$ with $m>rs$, choose a rational polynomial $Q$ of
degree at most $s$ and set
\[
          f(X)=Q(X)^r+\prod_{j=1}^{m}(X-j).
\]
It has degree $m$. The $m$ distinct points $(j,Q(j))$ lie on
$y^r=f(x)$ and on the feature hyperplane $y-Q(x)=0$. Thus a
smaller uniform bound cannot hold for this family as stated.
\end{remark}

\paragraph{The quintic instance.}
For any quintic $f$, take $r=2$, $s=2$. Then
\[
                  \Phi_2(x,y)=(1,x,x^2,y),
\]
and every proper subspace contains feature images of at most five
distinct affine rational points of $y^2=f(x)$. This includes the
monic squarefree quintic family used in the repository's Faltings
report \cite{repo-faltings}, but the proof does not need squarefreeness.
For a sharp numerical example take
\[
 f(X)=X^5-14X^4+85X^3-223X^2+274X-119
      =(X^2+1)^2+\prod_{j=1}^{5}(X-j).
\]
The points $(1,2),(2,5),(3,10),(4,17),(5,26)$ all lie in the
hyperplane $y-x^2-1=0$. The companion exact-arithmetic program
constructs an annihilator by minors and verifies the polynomial
identities for this example; the general proof is the argument above.

\subsection{Tuple-local covers and the remaining arithmetic theorem}

Let $\mathcal A$ be any additional decidable condition on the
points of $\mathcal C_{f,r}$. Examples might impose integrality,
congruence restrictions, or a particular approximation inequality.
We make no automatic assertion that these conditions yield a
small-subspace cover.

\begin{theorem}[The explicit finite cover interface]
\label{thm:superelliptic-cover}
Suppose that, for every finite list $L$ of distinct
$\mathcal C_{f,r}$-points satisfying $\mathcal A$, the feature
vectors $\Phi_s(L)$ admit a supplied cover by at most $M$ proper
rational linear subspaces. Then every such list has length at most
$mM$. It is enough to assume the cover assertion only for lists
of length $mM+1$. The covering subspaces may depend on the entire
input list; no global collection of exceptional subspaces is required.
\end{theorem}
\begin{proof}
For a supplied cover replace each subspace by the hyperplane from
Lemma~\ref{lem:section-annihilator}. Assign each point to the first
covering hyperplane by a bounded scan. The covering premise ensures
that an index is found. Each assigned sublist is distinct and
satisfies the corresponding hyperplane equation, so
Theorem~\ref{thm:superelliptic-section} bounds its length by $m$.
Summing over at most $M$ sublists gives $|L|\le mM$.
If the premise was stated only for length $mM+1$, any longer list
has a prefix of that length. The same argument would give
$mM+1\le mM$, a contradiction. The sum, assignment, and prefix
operations are bounded finite algorithms.
\end{proof}

The cover premise itself can be presented as a bounded arithmetic
predicate, without arbitrary subspace codes. For a list of $t$ feature
vectors, enumerate every subset of those vectors, compute its span,
and retain subsets of rank less than $s+2$. A span of vectors in a
proper supplied subspace remains proper: a nonzero annihilator of that
subspace also annihilates every linear combination of these vectors. A proper subspace
covering some of the tuple can be replaced by the span of the
points assigned to it. Thus one only needs $M$ subset masks, each
less than $2^t$, allowing the empty subset for padding. At most
$2^{tM}$ tuples of masks need be tested. All ranks and membership
conditions are decidable finite linear algebra. Equivalently,
replace each candidate span by its explicitly computed annihilator
hyperplane; this can only enlarge a covering set and remains proper.
This gives a fixed primitive recursive test
\[
       \operatorname{Cover}_{M,s}(L)\in\{0,1\}.
\]
The new theorem proves, in arithmetic,
\[
 \operatorname{ValidCurveList}_{f,r}(L)\ \wedge\
 \operatorname{Nodup}(L)\ \wedge\
 \operatorname{Cover}_{M,s}(L)=1
       \quad\Longrightarrow\quad |L|\le mM.
\]
It does not prove that the cover test succeeds for all lists with
$\mathcal A$. Proving that assertion is exactly where the requisite
quantitative approximation theorem must enter. This bounded
formulation identifies a precise interface for that future work.

If an explicitly listed exceptional set has at most $E$ points,
and the cover assertion applies outside it, the identical assignment
argument gives $E+mM$. For finitely many disjointly assigned
conditions with respective cover bounds $M_i$, the bound is
$E+m\sum_i M_i$. These are cardinality conclusions. They contain
no assertion locating all the exceptional points.

\subsection{Why a cardinality or gap bound still supplies no cutoff}

The repository correctly separates cardinality from a positive
height bound. Here is the exact general obstruction, including an
often tempting extra assumption. Let $T(e,n)$ mean that the fixed
universal machine with program/input code $e$ halts for the first
time at step $n$. Bounded simulation makes $T$ a fixed primitive
recursive predicate. For each $e$ it has at most one solution.
Consequently it satisfies every pairwise gap condition of the form
\[
 T(e,n)\wedge T(e,n')\wedge n<n'
             \quad\Longrightarrow\quad n'>G(e,n),
\]
for any proposed primitive recursive growth function $G$: there
are no two distinct solutions.

If a computable $B(e)$ bounded every solution of $T(e,n)$, simulate
the machine for $B(e)$ steps. This would decide whether it ever
halts. Such a decision procedure is impossible: a program which
halts exactly when that procedure predicts it will not halt,
applied to its own code, yields an immediate contradiction.
Therefore cardinality bounds, even together with arbitrarily strong
pairwise gap information, do not admit a general computable
conversion into height cutoffs. This argument uses one fixed bounded
simulator; it does not postulate a primitive recursive evaluator
for all primitive recursive functions.

Nothing here shows that a particular Diophantine cutoff is
unprovable. It shows why the finite-cover theorem must retain its
cardinality conclusion. Additional theorem-specific mathematics is
required to upgrade that conclusion
to a complete list or a uniform height bound.

\subsection{Lean definitions and proof obligations}

The following definitions suffice for a complete formalization path.
At the semantic level use \texttt{Rat}, finite coefficient arrays,
and finite lists of pairs. At the object level replace these by
the natural-number codes already described and use the same
bounded algorithms. The semantic and object-theoretic obligations
are distinct.

\begin{enumerate}
\item Define \texttt{QPoly} as a coefficient list, with
\texttt{coeff}, \texttt{eval}, \texttt{trim}, \texttt{degree},
\texttt{add}, \texttt{mul}, and \texttt{pow}. Zero is tested by a
Boolean fold over the coefficients. Implement
\texttt{syntheticDiv (p : QPoly) (z : Rat)} by the reverse coefficient
fold used in Lemma~\ref{lem:section-root-count}. Prove its coefficient
identity, the evaluation identity, and preservation of the leading
coefficient. Then prove \texttt{distinctRoots\_length\_le\_degree}
by iterating this operation over an input list, with the explicit
bounded invariant in that lemma.
\item Define rational matrices as dimensions plus flat arrays.
Implement minors, permutation determinants, bounded maximum-minor
rank, and \texttt{annihilator}. Prove determinant cofactor expansion
by finite sums, then the two identities used in
Lemma~\ref{lem:section-annihilator}. The target specification is
\texttt{rank W < n -> nonzero (annihilator W) = true}
together with \texttt{mulVec W (annihilator W) = zeroVec}.
The input $n>0$ is explicit.
\item Define
the curve, feature, and annihilator predicates by the finite operations
just described:
\begin{lstlisting}
onCurve f r (x,y) := (y^r == eval f x)
features s (x,y) := [1,x,...,x^s,y]
annihilates lam p := (dot lam (features s p) == 0)
\end{lstlisting}
The main specification quantifies over coefficient arrays and
finite lists:
if $r>0$, $\deg f=m>rs$, $\lambda$ has length $s+2$ and is
nonzero, and a duplicate-free list passes both Boolean tests, then
its length is at most $m$. Prove it by the decidable final-coefficient
split. In the nonzero case construct $(-P)^r-a^rf$; in the zero
case use a stable grouping operation with a proved partition-count
identity. No algebraic closure, real numbers, or abstract
finite-dimensional choice theorem is required.
\item Define \texttt{coverTest M s L} by the finite mask enumeration
above. Prove soundness by extracting the successful masks and their
annihilators. Prove completeness relative to supplied proper
subspaces by assigning each input point and replacing each group
by its span. Combine soundness with the section theorem to obtain
\texttt{coverTest\_implies\_length\_le}. A future quantitative
approximation development must separately prove that this Boolean
test is one on the required valid tuples.
\item Reify the rational arithmetic, coefficient folds, matrix
operations, and list grouping as particular primitive recursive
definitions. For every loop record its Boolean invariant and
compile its induction into a quantifier-free induction instance.
The human proofs above identify those instances. A compiled
semantic Lean theorem alone would not certify this final
object-theory translation.
\end{enumerate}

The root-count theorem, annihilator construction, superelliptic
section theorem, and finite-cover implication are proved
mathematically here. Their complete embedded Lean derivations are
formalization work. The uniform success of the quantitative
cover test is a further mathematical hypothesis, not part of that
formalization work and not claimed by this package.
